

\section{GIỚI THIỆU}

RISC (Reduced Instruction Set Computer) là một phương pháp thiết kế bộ xử lý hiện đại, 
tập trung vào việc giảm thiểu tập lệnh nhằm tăng tốc độ thực thi. 
Trong đề tài này, nhóm thực hiện thiết kế một RISC CPU đơn giản với 3-bit opcode và 5-bit toán hạng, 
tương ứng với 8 loại câu lệnh và 32 không gian địa chỉ.

\subsection{Tổng quan đề tài}
Trong bối cảnh ngành công nghiệp vi mạch và hệ thống nhúng đang phát triển mạnh mẽ, 
việc nắm vững nguyên lý thiết kế bộ xử lý từ cấp độ phần cứng là một nền tảng quan trọng 
đối với mỗi kỹ sư vi mạch và hệ thống nhúng. 
Thực hiện đề tài "Thiết kế RISC CPU đơn giản" giúp mô phỏng, kiểm chứng và minh họa toàn bộ chu kì làm việc của một CPU từ giai đoạn nạp lệnh (fetch), giải mã (decode), 
thực thi (execute) đến ghi kết quả (write-back).

RISC CPU được thiết kế sử dụng định dạng lệnh 8-bit bao gồm:
\begin{itemize}
    \item \textbf{3-bit opcode:} cho phép mã hóa tối đa 8 lệnh khác nhau.
    \item \textbf{5-bit toán hạng:} cung cấp 32 địa chỉ cho bộ nhớ và thanh ghi.
\end{itemize}

Kiến trúc loại bỏ các lệnh phức tạp, giữ lại các lệnh cơ bản giúp các khối chức năng được thiết kế tinh gọn, đơn giản hơn và tăng tốc quá trình thực thi lệnh.
\subsection{Mục tiêu của nhóm}
Nhóm đặt ra các mục tiêu cụ thể sau cho đề tài:
\begin{enumerate}
    \item \textbf{Nắm vững cơ sở lý thuyết}: Hiểu được nguyên lý hoạt động của bộ xử lý theo thiết kế RISC 
    làm cơ sở để xây dựng các sơ đồ khối.
    \item \textbf{Thiết kế tập lệnh (ISA)}: Xây dựng một tập lệnh gồm 8 lệnh cơ bản, bao gồm các lệnh số học, logic, truy xuất bộ nhớ và điều khiển luồng chương trình.
    \item \textbf{Hiện thực thiết kế bằng HDL}: Sử dụng ngôn ngữ Verilog để xây dựng hệ thống.
    \item \textbf{Mô phỏng và kiểm thử}: Viết các testbench để kiểm tra từng khối chức năng cũng như toàn bộ hệ thống.
    \item \textbf{Trình bày kết quả rõ ràng}: Viết báo cáo bao gồm đầy đủ các sơ đồ, giản đồ sóng và giải thích nguyên lý.
\end{enumerate}

% \subsection{Lý do chọn đề tài}
% Nhóm lựa chọn đề tài thiết kế RISC CPU xuất phát từ ba lý do chính:

% \textbf{Thứ nhất, giá trị học thuật cao.} Bộ xử lý là trái tim của mọi hệ thống máy tính. Việc tự tay thiết kế một CPU — dù ở quy mô nhỏ — giúp chúng tôi hiểu thực sự cách phần cứng và phần mềm tương tác với nhau, thay vì chỉ tiếp cận lý thuyết qua sách vở.

% \textbf{Thứ hai, tính thực tiễn và ứng dụng.} Kiến trúc RISC là nền tảng của các bộ xử lý phổ biến nhất hiện nay như ARM (dùng trong smartphone), RISC-V (đang được áp dụng rộng rãi trong nghiên cứu và công nghiệp). Nắm vững RISC là bước khởi đầu thiết thực để tiếp cận các kiến trúc thực tế này.

% \textbf{Thứ ba, phạm vi phù hợp với thời gian và nguồn lực.} Với định dạng lệnh 8-bit và tập lệnh giới hạn ở 8 lệnh, đề tài có độ phức tạp vừa đủ — đủ thách thức để học hỏi sâu, nhưng không quá lớn để mất kiểm soát trong khuôn khổ học phần. Điều này cho phép nhóm tập trung vào chất lượng thiết kế thay vì chạy theo số lượng tính năng.

% Với những định hướng trên, nhóm kỳ vọng đề tài không chỉ là một bài tập kỹ thuật hoàn chỉnh, mà còn là nền tảng vững chắc để mỗi thành viên tự tin tiếp cận các hệ thống vi xử lý phức tạp hơn trong tương lai.

\subsection{Các phần mềm và công cụ sử dụng}
\begin{itemize}
    \item \textbf{Xilinx Vivado Design Suite:} Phần mềm để soạn thảo code Verilog, đồng thời hỗ trợ biên dịch và chạy mô phỏng
    \item \textbf{Draw.io: } Công cụ vẽ sơ đồ khối và kết nối các khối chức năng trong một RISC CPU.
\end{itemize}

\subsection{Thiết bị phần cứng}
\textbf{Bo mạch FPGA Digilent Arty Z7: }Bo mạch được sử dụng để nạp và kiểm thử thiết kế RISC CPU trực tiếp trên phần cứng thực tế.

\subsection{Các chức năng của sản phẩm}
Hệ thống RISC CPU được thiết kế bao gồm các chức năng cốt lõi sau:
\begin{itemize}
    \item \textbf{Nạp lệnh (Instruction Fetch):} Lấy lệnh tiếp theo từ bộ nhớ.
    \item \textbf{Giải mã lệnh (Instruction Decode):} Phân tích mã lệnh để xác định thao tác cần thực hiện.
    \item \textbf{Lấy dữ liệu toán hạng (Operand Fetch):} Truy xuất dữ liệu từ bộ nhớ (nếu cần).
    \item \textbf{Thực thi câu lệnh (Execute):} Xử lý các phép toán thông qua khối ALU.
    \item \textbf{Lưu trữ kết quả (Store):} Ghi dữ liệu về bộ nhớ hoặc Accumulator.
\end{itemize}

% \vspace{0.3cm}
% \noindent \textbf{Tập lệnh hỗ trợ:} Hệ thống được thiết kế để tự động dừng chương trình khi gặp lệnh HALT và hỗ trợ 8 lệnh cơ bản như bảng sau:

% \begin{table}[h!]
%     \centering
%     \renewcommand{\arraystretch}{1.3}
%     \begin{tabular}{|c|c|l|}
%         \hline
%         \textbf{STT} & \textbf{Lệnh} & \textbf{Chức năng} \\ \hline
%         1 & \texttt{HLT} & Dừng chương trình \\ \hline
%         2 & \texttt{SKZ} & Bỏ qua lệnh tiếp theo nếu Accumulator = 0 \\ \hline
%         3 & \texttt{ADD} & Cộng dữ liệu \\ \hline
%         4 & \texttt{AND} & Phép toán AND logic \\ \hline
%         5 & \texttt{XOR} & Phép toán XOR logic \\ \hline
%         6 & \texttt{LDA} & Nạp dữ liệu vào Accumulator \\ \hline
%         7 & \texttt{STO} & Lưu dữ liệu từ Accumulator ra bộ nhớ \\ \hline
%         8 & \texttt{JMP} & Nhảy đến địa chỉ chỉ định \\ \hline
%     \end{tabular}
%     \caption{Tập lệnh của RISC CPU}
% \end{table}

